% Part: computability
% Chapter: computability-theory
% Section: ce-closed-cup-cap

\documentclass[../../../include/open-logic-section]{subfiles}

\begin{document}

\olfileid{cmp}{thy}{clo}

\olsection[CE سیٹاں دا اتحاد تے اشتراک]{قابلِ حساب گنتی وار
  فہرست بنن جوگ سیٹ اتحاد تے اشتراک تحت بند نیں}

اگلا قضیہ قابلِ حساب گنتی وار فہرست بنن جوگ
سیٹاں دے مجموعے دیاں کجھ بندش دیاں خاصیتاں
دسدا اے۔

\begin{thm}
من لو $A$ تے $B$ قابلِ حساب گنتی وار
فہرست بنن جوگ نیں۔ تاں $A \cap B$
تے $A \cup B$ وی ایداں دے نیں۔
\end{thm}

\begin{proof}
\olref[eqc]{thm:ce-equiv} سانوں قابلِ حساب گنتی وار
فہرست بنن جوگ سیٹاں دیاں وکھ وکھ تعریفاں
ورتن دی اجازت دیندا اے۔ مثال دے طور اُتے
اسی کئی وکھرے ثبوت دیاں گے۔
% PNB-CMP-017: the first two source constructions assume nonempty input sets.
جے دوہاں وچوں کوئی اک سیٹ خالی ہووے تاں
اوہناں دا اتحاد دوجا سیٹ اے تے اشتراک خالی؛
ایس صورت وچ دعویٰ فوراً نکلدا اے۔ پہلے
دو ثبوتاں دیاں تعمیرات وچ اس لئی غیر خالی
صورت نوں ویکھنا کافی اے۔

پہلے ثبوت لئی من لو $A$ دی گنتی وار فہرست
قابلِ حساب فنکشن~$f$ بناندا اے، تے $B$
دی گنتی وار فہرست قابلِ حساب فنکشن~$g$
بناندا اے۔ ایہ رکھو:
\begin{align*}
h(x) & = \umin{y}{(f(y) = x \lor g(y) = x)} \text{ تے}\\
j(x) & = \umin{y}{(f((y)_0) = x \land g((y)_1) = x)}.
\end{align*}
تاں $A \cup B$، $h$ دا دائرۂ تعریف
اے، تے $A \cap B$، فنکشن~$j$
دا دائرۂ تعریف اے۔

\begin{explain}
حسابی لحاظ نال ایہ ہو رہیا اے: جے ساڈے کول
$A$ تے $B$ دیاں گنتی وار فہرستاں بناؤن والے
طریقے ہون، تاں کسے رکن $x$ دے
$A \cup B$ وچ ہون بارے نیم فیصلہ
کرن لئی دوناں وچوں کسے وی فہرست وچ $x$
لبھ سکدے آں؛ تے $x$ دے $A \cap B$
وچ ہون بارے نیم فیصلہ کرن لئی دوناں
فہرستاں وچ اکّو ویلے $x$ لبھ سکدے آں۔
\end{explain}

دوجے ثبوت لئی فیر من لو $A$ دی
فہرست~$f$ بناندا اے تے $B$ دی
فہرست~$g$۔ ایہ رکھو:
\[
k(x) = \begin{cases}
f(x/2) & \text{جے $x$ جفت ہووے} \\
g((x-1)/2) & \text{جے $x$ طاق ہووے۔}
\end{cases}
\]
تاں $k$، $A \cup B$ دی فہرست
بناندا اے؛ خیال ایہ اے کہ $k$،
$f$ تے~$g$ دیاں دتیاں فہرستاں
وارو واری ورتدا اے۔ $A \cap
B$ دی فہرست بنانا ذرا اوکھا اے۔
جے $A \cap B$ خالی اے، تاں ایہ
تعریف مطابق ای گنتی وار فہرست بنن جوگ
اے۔ نہیں تاں $c$ نوں $A \cap B$
دا کوئی وی رکن چن لو، تے $l$ دی
تعریف ایویں کرو:
\[
l(x) = \begin{cases}
f((x)_0) & \text{جے $f((x)_0) = g((x)_1)$} \\
c & \text{ہور صورت وچ۔}
\end{cases}
\]
حسابی لحاظ نال $l$، $f$ تے $g$
دیاں فہرستاں دے رکناں دیاں جوڑیاں
ویکھدا اے تے جتھے دوویں ملن، اوہ
قدر واپس کردا اے؛ نہیں تاں صرف $c$
واپس کردا اے، یعنی اوہی مقرر قدر دُہراندا اے۔

آخری ثبوت لئی من لو $A$، جزوی
فنکشن $m(x)$ دا \emph{دائرۂ تعریف}
اے تے $B$، جزوی فنکشن $n(x)$
دا دائرۂ تعریف اے۔ تاں $A \cap B$،
جزوی فنکشن $m(x) +
n(x)$ دا دائرۂ تعریف اے۔

\begin{explain}
حسابی لحاظ نال جے $A$ اوہ قدراں
دا سیٹ اے جنہاں اُتے $m$ رُکدا اے،
تے $B$ اوہ قدراں دا سیٹ اے جنہاں
اُتے $n$ رُکدا اے، تاں $A \cap B$
اوہ قدراں دا سیٹ اے جنہاں اُتے
دونوں طریقے رُکدے نیں۔
\end{explain}

$A \cup B$ نوں رُکّن والیاں قدراں
دے سیٹ دے روپ وچ لکھنا ودھ اوکھا اے،
کیوں جے $m$ تے $n$ دے حسابی عملاں دی
متوازی طور اُتے مشینی نقل کرنی پیندی اے۔ من لو
$d$، $m$ دا اشاریہ اے تے $e$،
$n$ دا اشاریہ اے؛ ہور لفظاں وچ،
$m = \cfind{d}$ تے $n = \cfind{e}$۔
تاں $A \cup B$ اِس فنکشن دا دائرۂ
تعریف اے:
\[
p(x) = \umin{y}{(T(d,x,y) \lor T(e,x,y))}.
\]
\begin{explain}
حسابی لحاظ نال، انپٹ $x$ اُتے $p$،
یا $m$ لئی یا $n$ لئی رُک جان والی
حسابی کارروائی لبھدا اے، تے دوہاں
وچوں کوئی وی لبھ جائے تاں رُک جاندا اے۔
\end{explain}
\end{proof}

\end{document}
